\documentclass{article}
\usepackage{amsmath, amssymb, amsthm}
\usepackage{hyperref}
\usepackage{geometry}
\geometry{margin=1in}

\title{Erd\H{o}s Problem \#151}
\date{Last updated: 2025-08-31}
\author{Source: erdosproblems.com}

\begin{document}
\maketitle

\noindent\textbf{Status:} OPEN \quad \textbf{No prize}

\noindent\textbf{Tags:} graph theory

\medskip

\noindent\textbf{Problem Statement:}

For a graph $G$ let $\tau(G)$ denote the minimal number of vertices that include at least one from each maximal clique of $G$ on at least two vertices (sometimes called the clique transversal number). Let $H(n)$ be maximal such that every triangle-free graph on $n$ vertices contains an independent set on $H(n)$ vertices. If $G$ is a graph on $n$ vertices then is\[\tau(G)\leq n-H(n)?\]




It is easy to see that $\tau(G) \leq n-\sqrt{n}$. Note also that if $G$ is triangle-free then trivially $\tau(G)\leq n-H(n)$. This is listed in \cite{Er88} as a problem of Erd\H{o}s and Gallai, who were unable to make progress even assuming $G$ is $K_4$-free. There Erd\H{o}s remarked that this conjecture is 'perhaps completely wrongheaded'. It later appeared as Problem 1 in \cite{EGT92}. The general behaviour of $\tau(G)$ is the subject of [610] .




References

[EGT92] Erd\H{o}s, Paul and Gallai, Tibor and Tuza, Zsolt, Covering the cliques of a graph with vertices . Discrete Math. (1992), 279-289.

[Er88] Erd\H{o}s, P, Problems and results in combinatorial analysis and graph theory . Discrete Math. (1988), 81-92.

\noindent\textbf{Related OEIS sequences:} possible


\bigskip
\noindent\small{Source: \url{https://www.erdosproblems.com/151}}
\end{document}
